% ********** Rozdział 2 **********
\chapter{Użyte narzędzia i~kod źródłowy}\label{ch:narzedzia}

\section{Wykaz narzędzi}

Program korzysta wyłącznie z~platformy .NET i~jej biblioteki standardowej. Nie wymaga żadnych dodatkowych pakietów NuGet ani serwera bazy danych. Tabela \ref{tab:narzedzia} zawiera pełny wykaz narzędzi.

\begin{table}[H]
	\begin{center}
		\caption{\footnotesize Wykaz użytych narzędzi}
		\label{tab:narzedzia}
		\begin{small}
			\begin{tabular}{|l|p{10.5cm}|} \hline
				\textbf{Narzędzie} & \textbf{Zastosowanie} \\ \hline
				.NET 9 SDK, język C\# & Kompilacja i~uruchamianie programu. \\ \hline
				System.Text.Json & Zapis i~odczyt danych w~formacie JSON; biblioteka wchodzi w~skład platformy .NET. \\ \hline
				Plik JSON & Trwałe przechowywanie danych; baza danych nie jest używana. \\ \hline
				Git, GitHub & Kontrola wersji i~udostępnienie kodu źródłowego. \\ \hline
				LaTeX (pdfTeX) & Przygotowanie tej dokumentacji na podstawie szablonu uczelnianego. \\ \hline
				Python (Pillow, numpy), Google Chrome & Automatyczne wykonanie zrzutów ekranu do dokumentacji (skrypt \texttt{dokumentacja/narzedzia/zrzuty.py}, opis w~pliku \texttt{ZRZUTY.md}). Nie są potrzebne do działania programu. \\ \hline
			\end{tabular}
		\end{small}
	\end{center}
\end{table}

\section{Repozytorium z~kodem źródłowym}

Kod źródłowy projektu wraz ze źródłami tej dokumentacji znajduje się w~repozytorium:

\begin{center}
	\url{\repozytorium}
\end{center}

\section{Uruchomienie programu}

Do uruchomienia programu potrzebny jest zainstalowany pakiet .NET 9 SDK. Po pobraniu repozytorium program uruchamia się z~jego głównego katalogu poleceniami:

\begin{lstlisting}
git clone https://github.com/BobnikTheIII/rejestr-urzadzen-sieciowych.git
cd rejestr-urzadzen-sieciowych
dotnet run --project src/RejestrUrzadzen
\end{lstlisting}

Program można też otworzyć w~środowisku Visual Studio za pomocą pliku \texttt{RejestrUrzadzen.sln}.

\subsection{Położenie pliku danych}

Domyślnie dane są zapisywane w~pliku \texttt{urzadzenia.json} w~katalogu danych aplikacji bieżącego użytkownika:

\begin{itemize}
	\item Windows: \texttt{\%LOCALAPPDATA\%\textbackslash RejestrUrzadzen\textbackslash urzadzenia.json},
	\item Linux: \texttt{\textasciitilde/.local/share/RejestrUrzadzen/urzadzenia.json}.
\end{itemize}

Inny plik można wskazać jako argument programu. Jeśli plik nie istnieje, zostanie utworzony przy pierwszym zapisie:

\begin{lstlisting}
dotnet run --project src/RejestrUrzadzen -- moj_rejestr.json
\end{lstlisting}

Ścieżka używanego pliku jest zawsze widoczna w~menu głównym.

% ********** Koniec rozdziału **********
